

\documentclass[../thesis.tex]{subfiles}
\begin{document}
seL4, the "world's fastest microkernel" \citep{Heiser_20:sel4wp} and the first formally verified microkernel \citep{Klein_AEHCDEEKNSTW_10} has shown great applicability as a base for an Operating System for security and reliability critical systems in the embedded/IoT space through its Operating System Implementation LionsOS \citep{Klein_AKMHF_18,Heiser_VCJGNLDZAZWPDB_25}.  
% talk about how microkernel use is limited in other spaces, general etc.
Despite early prevalence and the implication that microkernels were to be the future for their theoretically superior design, in the modern day, pure microkernel use is relatively obscure, reserved for mostly embedded use cases.
% why is this: probably due to performance and portability
Following the Mach 3.0 debacle \citep{Chen_Bershad_93}, the public perception of microkernels shifted, this paper attributing the poor performance of Mach in comparison to Ultrix due to Microkernel design, namely \gls{ipc} and the distributed nature of Mach 3.0 which was redesigned from a hybrid kernel to a microkernel, leading to more cache misses and higher instruction counts compared to Ultrix for the same tasks. While it is true that \gls{ipc} is slower than the equivalent mode switch for a monolithic kernel, Liedtke demonstrated through his L3 microkernel that \gls{ipc} can be improved by an order of magnitude \citep{Liedtke_93} and that Mach 3.0's approach to microkernels was incorrect, attributing the user-kernel cache competition to a large microkernel \citep{Liedtke_96b}. This opened up the possibility that the \gls{ipc} performance differences may be amortised by other more significant processes, but it seems that this idea has not been extensively explored and the idea that microkernels being less performant remained.

% ai and security.

Due to the inherent vulnerabilities of monolithic kernels and its non-modular design, there have been movements towards microkernels even in the general purpose operating system space. Windows NT and XNU are both hybrid kernels and more recently the HongMeng kernel claims to be a general purpose microkernel which shows comparable performance to Linux in its benchmarks. \citep{Chen_JWLLLWHLYWYPX_24}. However, it is clear that HongMeng forgoes many microkernel principles with the purpose of achieving higher performance, making design choices that go against the principle of minimality, increasing the size of the code base to ~90k\gls{sloc} and therefore reducing the feasibility of verification and more importantly increasing the probability of bugs being present in the kernel. Their approach appears to contradict Liedtke's findings, despite their good performance unlike Mach 3.0.

% the LIONSOS paper showed potential about the performance of seL4 as a kernel for performance critical tasks
The question arises: \textit{Is it possible to create a performant microkernel-based general purpose operating system without compromising on security?} Recently LionsOS showed that in the context of a static operating system that a microkernel based system is able to outperform Linux at least in network benchmarks resulting in 100,000s context switches per second, at least as challenging as the 60,000 syscalls in the HongMeng paper \citep{Heiser_VCJGNLDZAZWPDB_25,Chen_JWLLLWHLYWYPX_24}. LionsOS' superior benchmarks in almost all cases in what is possibly a worst case scenario for a microkernel-based system raises doubt about the significance of IPC regarding performance as is currently widely believed. However, this is not a perfect comparison as LionsOS is a static system and Linux a dynamic system.
% talk about how there is potential for microkernels to be even faster than monolithic.
% this thesis is aimed at attempting to prove that microkernels are not only suitable but furthermore better than monolithic kernels even for performance critical workloads

The aim of this thesis is to prove that microkernels can be at least as performant as monolithic systems even as a general purpose operating system without compromising on security. This will be done through proof of concept implementations of operating system components on top of an unmodified seL4 kernel and benchmarking them to prove that microkernels can indeed be performant.
% We prove this by debunking the negative claims made about microkernels which may/may not have legitimate backing
% we will also do this without major modifications of the seL4 kernel - adhering to the minimality principle

\end{document}